% Shared assertions for the star (safe-sizing) routes of every enclosure.
%
% A fixture inputs this file in its preamble and calls \KKSfitsuite in the
% document body.  The suite drives the *public* commands, records what the
% package actually did to their material, and compares that with an interior
% measured from a rendered outline, so that neither the fit code's own limits
% nor a constant copied out of KKsymbols.sty can make a check pass by itself.
%
% How the interior is measured.  PGF adds the stroke to a picture's bounding
% box, and `outer sep' -- which it initialises to half the line width -- is
% added to a node's own box.  Drawing the same node three times therefore
% yields everything the geometry needs without asking the package:
%
%   draw=none, outer sep=0pt : the shape's box, i.e. its half-extent
%   draw,      outer sep=0pt : the same plus half the stroke on either side
%   draw                     : the background path plus half the stroke
%
% The first two differ by exactly the stroke that is really used, and taking
% half of that back off the third gives the path that is really drawn.  The
% usable interior is that path taken in by half the stroke.
\makeatletter
\newsavebox{\KKSfitprobe}
\newsavebox{\KKSfitprobedraw}
\newsavebox{\KKSfitprobepath}
\newsavebox{\KKSfitout}
\newlength{\KKSfitstroke}
\newlength{\KKSfitpath}
\newlength{\KKSfittmp}
\newlength{\KKSfitgotw}
\newlength{\KKSfitgoth}
\newlength{\KKSfitnatw}
\newlength{\KKSfitnath}
\newlength{\KKSfitlima}
\newlength{\KKSfitlimrc}
\newlength{\KKSfitlimcc}
\newlength{\KKSfitlimband}
\newlength{\KKSfitlimwide}
\newif\ifKKSfitwideheld
\def\KKSfittaken{0}
% How far the recorded box reaches on the axes rather than at its
% corner.  This is the test the ink guard applies, and it is the one a
% shape applies to a lone glyph, whose ink stops short of its corners.
\newcommand{\KKSfitsetreachaxis}[1]{%
  \setlength{\KKSfittmp}{.5\KKSfitrw}%
  \ifdim.5\KKSfitrh>\KKSfittmp\setlength{\KKSfittmp}{.5\KKSfitrh}\fi
  \KKSfitsetratio{\KKSfitreach}{\KKSfittmp}{#1}}
\newlength{\KKSfitbandreach}
\newlength{\KKSfitinner}
\newlength{\KKSfitcap}
\newlength{\KKSfitusedlw}
\newlength{\KKSfitdeclw}
\newsavebox{\KKSfitfinal}
\newlength{\KKSfitfinw}
\newlength{\KKSfitfinh}
\newlength{\KKSfitrw}
\newlength{\KKSfitrh}
\newlength{\KKSfitaspect}
\newlength{\KKSfitgotaspect}
\newlength{\KKSfitnataspect}
\newlength{\KKSfitfinaspect}
\newlength{\KKSfitratio}
\newlength{\KKSfitshrunk}
\newcommand{\KKSfitfail}[2]{\PackageError{\jobname}{Star fit `#1': #2}{}}
% pgf rounds a circle's radius through low-precision integer arithmetic, and
% every dimension here is rounded to scaled points twice over, so comparisons
% carry a tolerance far below anything that could be seen.
\newcommand{\KKSfitslack}{0.02pt}
% The tolerance for the two-sided comparison of the limit a command aimed at
% with the one its own outline really offers.  These are the same quantity
% reached by two different routes -- pgf's shape arithmetic and the package's
% closed form -- and over the fixtures here the two never disagree by as much
% as half a thousandth of a point, so ten times that is left for a font or a
% frame width that rounds less kindly.  It has to stay this tight: the whole
% point of the comparison is that it is far smaller than the shift a swapped
% line-width key or a swapped corner model causes.
\newcommand{\KKSfitexactslack}{0.005pt}

% ---------------------------------------------------------------------------
% Record what the package really applied.  \kks@fit@emit is the single point
% every star route funnels through, so wrapping it captures the natural size,
% the limits the command computed, and the box the reader finally sees.
% ---------------------------------------------------------------------------
% The line width the command resolved for its own setup.  A wrapper that
% reads another command's key is invisible in the finished page whenever the
% two keys happen to agree, so it is recorded here and compared with the width
% the drawing node is declared to use.
\let\KKSfit@linewidth\kks@fit@linewidth
\renewcommand{\kks@fit@linewidth}[1]{%
  \KKSfit@linewidth{#1}%
  \global\KKSfitusedlw=\kks@fit@lw
  \global\let\KKSfitlwseen\@empty}

% The final, public size of the material.  \kks@fit@emit sees only what the
% uniform fit produced; the shrink-only visual guard may take that box in
% further before the reader gets it, so the guarded routes are wrapped too and
% the box they really emit is measured.
\NewCommandCopy\KKSfit@guardcircle\kks@visual@guard@circle
\RenewDocumentCommand{\kks@visual@guard@circle}{ m m m m m m }{%
  \sbox{\KKSfitfinal}{\KKSfit@guardcircle{#1}{#2}{#3}{#4}{#5}{#6}}%
  \global\KKSfitfinw=\wd\KKSfitfinal
  \global\KKSfitfinh=\dimexpr\ht\KKSfitfinal+\dp\KKSfitfinal\relax
  \usebox{\KKSfitfinal}}
\NewCommandCopy\KKSfit@guardbracket\kks@visual@guard@bracket
\RenewDocumentCommand{\kks@visual@guard@bracket}{ m m m m m m }{%
  \sbox{\KKSfitfinal}{\KKSfit@guardbracket{#1}{#2}{#3}{#4}{#5}{#6}}%
  \global\KKSfitfinw=\wd\KKSfitfinal
  \global\KKSfitfinh=\dimexpr\ht\KKSfitfinal+\dp\KKSfitfinal\relax
  \usebox{\KKSfitfinal}}

\let\KKSfit@emit\kks@fit@emit
\renewcommand{\kks@fit@emit}[2]{%
  \sbox{\KKSfitout}{\KKSfit@emit{#1}{#2}}%
  \global\KKSfitgotw=\wd\KKSfitout
  \global\KKSfitgoth=\dimexpr\ht\KKSfitout+\dp\KKSfitout\relax
  \global\KKSfitnatw=\kks@fit@wd
  \global\KKSfitnath=\kks@fit@ht
  \global\KKSfitlima=\kks@fit@a
  \global\KKSfitlimrc=\kks@fit@rc
  \global\KKSfitlimcc=\kks@fit@cx
  % The band as the fit captured it, not the live register: a nested
  % enclosure has already overwritten the register by the time we get here.
  \global\KKSfitlimband=\kks@fit@bband\relax
  \global\KKSfitlimwide=\kks@fit@wwide\relax
  \global\let\KKSfittaken\kks@fit@widetaken
  \global\let\KKSfitseen\@empty
  \usebox{\KKSfitout}}

% ---------------------------------------------------------------------------
% Measure one enclosure's drawn outline.
%
%   #1 name          #2 kind (circle, rect or diamond)
%   #3 node options, exactly as the enclosure states them, but with sharp
%      corners on a diamond: pgf tracks a rounded path through its control
%      points and would report a tip the pen never reaches, so a rounded
%      diamond tip is accounted for by #6 instead
%   #4 side of the square text box the enclosure hands the node
%   #5 the node's own `scale' (the doubled circle carries one; pgf transforms
%      the path but not the stroke, so the interior is converted back into the
%      unscaled units the material is laid out in)
%   #6 corner radius of the drawn path, 0pt for a sharp corner
%   #7 the line width the drawing node is declared with, so that the width the
%      command resolved for its own setup can be compared with it
% ---------------------------------------------------------------------------
\newcommand{\KKSfitnode}[3]{%
  \tikz{\node[#1,#2] (probe) {%
    \vbox to #3{\vss\hbox to #3{\hss}\vss}};}}
\newcommand{\KKSfitshape}[7]{%
  \begingroup
    \pgfmathsetlength{\KKSfitinner}{#4}%
    \sbox{\KKSfitprobe}{\KKSfitnode{#3}{draw=none,outer sep=0pt}{\KKSfitinner}}%
    \sbox{\KKSfitprobedraw}{\KKSfitnode{#3}{draw,outer sep=0pt}{\KKSfitinner}}%
    \sbox{\KKSfitprobepath}{\KKSfitnode{#3}{draw}{\KKSfitinner}}%
    \setlength{\KKSfitstroke}{%
      \dimexpr\wd\KKSfitprobedraw-\wd\KKSfitprobe\relax}%
    \setlength{\KKSfitpath}{%
      \dimexpr.5\wd\KKSfitprobepath-.5\KKSfitstroke\relax}%
    \ifdim\KKSfitstroke<\z@\KKSfitfail{#1}{measured a negative stroke}\fi
    \ifdim\KKSfitpath<\z@\KKSfitfail{#1}{measured a negative outline}\fi
    % The interior: the drawn path taken in by half the stroke.  A slanted
    % diamond edge moves .707 of a whole stroke along its diagonal, and a
    % rounded tip costs a further .414 of the radius the inner edge keeps.
    \edef\KKSfitkind{#2}%
    \def\KKSfitkinddiamond{diamond}%
    \ifx\KKSfitkind\KKSfitkinddiamond
      \pgfmathsetlength{\KKSfittmp}{%
        (\KKSfitpath - 0.7071068*\KKSfitstroke
           - 0.4142136*max(0, (#6) - .5*\KKSfitstroke)) / (#5)}%
    \else
      \pgfmathsetlength{\KKSfittmp}{(\KKSfitpath - .5*\KKSfitstroke) / (#5)}%
    \fi
    \global\expandafter\let\csname KKSfit@kind@#1\endcsname\KKSfitkind
    \expandafter\xdef\csname KKSfit@usable@#1\endcsname{\the\KKSfittmp}%
    % Centre and radius of the drawn corner arc, in the same units.
    \pgfmathsetlength{\KKSfittmp}{max(0, (#6) - .5*\KKSfitstroke) / (#5)}%
    \expandafter\xdef\csname KKSfit@rc@#1\endcsname{\the\KKSfittmp}%
    \pgfmathsetlength{\KKSfittmp}{(\KKSfitpath - (#6)) / (#5)}%
    \expandafter\xdef\csname KKSfit@cc@#1\endcsname{\the\KKSfittmp}%
    % The historical limit: the largest half-extent of the shape's own metric
    % that the inner box the node is handed can stand for.  A square of side s
    % has half-diagonal s/sqrt(2) on a circle, half-extent s/2 in a rectangle,
    % and is inscribed exactly in a diamond of half-diagonal s.
    \ifx\KKSfitkind\KKSfitkinddiamond
      \setlength{\KKSfitcap}{\KKSfitinner}%
    \else
      \def\KKSfitkindcircle{circle}%
      \ifx\KKSfitkind\KKSfitkindcircle
        \pgfmathsetlength{\KKSfitcap}{0.5*sqrt(2)*\KKSfitinner}%
      \else
        \setlength{\KKSfitcap}{.5\KKSfitinner}%
      \fi
    \fi
    \expandafter\xdef\csname KKSfit@cap@#1\endcsname{\the\KKSfitcap}%
    \pgfmathsetlength{\KKSfittmp}{max(0, #7)}%
    \expandafter\xdef\csname KKSfit@lw@#1\endcsname{\the\KKSfittmp}%
    \expandafter\xdef\csname KKSfit@stroke@#1\endcsname{\the\KKSfitstroke}%
    \pgfmathsetlength{\KKSfittmp}{#6}%
    \expandafter\xdef\csname KKSfit@rcdecl@#1\endcsname{\the\KKSfittmp}%
    \typeout{KKS-FIT-REF:#1 kind=#2 stroke=\the\KKSfitstroke\space
      path=\the\KKSfitpath\space
      usable=\csname KKSfit@usable@#1\endcsname\space
      cap=\csname KKSfit@cap@#1\endcsname\space
      lw=\csname KKSfit@lw@#1\endcsname}%
  \endgroup}

% Every ratio in this file is formed here rather than by pgfmath.  pgfmath's
% division collapses on operands whose quotient is close to one -- it answers
% 3.57007/3.57011 with 1.1 -- and these assertions are made of exactly such
% quotients, a box measured against the room it was fitted into.  TeX's own
% scaling is exact instead: #1 is set to the quotient of #2 by #3 as a length
% whose value in points *is* that quotient, to within one scaled point.  A
% region with nothing left inside it can hold nothing, so it reads as an
% unbounded reach rather than as a division by zero.
\newlength{\KKSfitnum}
\newlength{\KKSfitden}
\newcommand{\KKSfitsetratio}[3]{%
  \setlength{\KKSfitnum}{#2}%
  \setlength{\KKSfitden}{#3}%
  \ifdim\KKSfitden>\z@
    \setlength{#1}{\numexpr\KKSfitnum*65536/\KKSfitden\relax sp}%
  \else
    \setlength{#1}{\maxdimen}%
  \fi}
\newcommand{\KKSfitabs}[1]{\ifdim#1<\z@\setlength{#1}{-#1}\fi}

% How far the recorded box, centred on the shape, gets towards the boundary of
% a region: 1pt is exactly on it.  A box on a circle of radius a reaches
% sqrt(w^2+h^2)/2a; on a diamond with half-diagonals a it reaches w/2a + h/2a;
% in a rounded rectangle the sides limit it unless its corner passes both arc
% centres, and then the arc does.  #1 kind, #2 half-extent, #3 arc centre,
% #4 arc radius.
\newlength{\KKSfitreach}
\newlength{\KKSfitcorner}
\newcommand{\KKSfitsetreach}[4]{%
  \edef\KKSfitkind{#1}%
  \def\KKSfitkinddiamond{diamond}%
  \def\KKSfitkindcircle{circle}%
  \ifx\KKSfitkind\KKSfitkinddiamond
    \KKSfitsetratio{\KKSfitreach}
      {\dimexpr.5\KKSfitrw+.5\KKSfitrh\relax}{#2}%
  \else
    \ifx\KKSfitkind\KKSfitkindcircle
      \pgfmathsetlength{\KKSfittmp}{%
        .5*sqrt(\KKSfitrw*\KKSfitrw + \KKSfitrh*\KKSfitrh)}%
      \KKSfitsetratio{\KKSfitreach}{\KKSfittmp}{#2}%
    \else
      \setlength{\KKSfittmp}{.5\KKSfitrw}%
      \ifdim.5\KKSfitrh>\KKSfittmp\setlength{\KKSfittmp}{.5\KKSfitrh}\fi
      \KKSfitsetratio{\KKSfitreach}{\KKSfittmp}{#2}%
      \ifdim.5\KKSfitrw>#3\relax
        \ifdim.5\KKSfitrh>#3\relax
          \pgfmathsetlength{\KKSfittmp}{%
            sqrt((.5*\KKSfitrw - (#3))*(.5*\KKSfitrw - (#3))
                 + (.5*\KKSfitrh - (#3))*(.5*\KKSfitrh - (#3)))}%
          \KKSfitsetratio{\KKSfitcorner}{\KKSfittmp}{#4}%
          \ifdim\KKSfitcorner>\KKSfitreach
            \setlength{\KKSfitreach}{\KKSfitcorner}%
          \fi
        \fi
      \fi
    \fi
  \fi}
% The box \KKSfitsetreach is to judge.
\newcommand{\KKSfitjudgefit}{%
  \setlength{\KKSfitrw}{\KKSfitgotw}\setlength{\KKSfitrh}{\KKSfitgoth}}
\newcommand{\KKSfitjudgefinal}{%
  \setlength{\KKSfitrw}{\KKSfitfinw}\setlength{\KKSfitrh}{\KKSfitfinh}}

% ---------------------------------------------------------------------------
% One check: run the public command on the material and judge the result.
%   #1 name (as declared below)   #2 the public command, star included
%   #3 the material               #4 the largest scale the route may apply
% ---------------------------------------------------------------------------
\newcommand{\KKSfitcheck}[4]{%
  \begingroup
    \global\let\KKSfitseen\relax
    \global\let\KKSfitlwseen\relax
    \global\KKSfitfinw=-\maxdimen
    \setbox\z@=\hbox{#2{#3}}%
    \ifx\KKSfitseen\relax
      \KKSfitfail{#1/#3}{the star route never reached the uniform fit}%
    \fi
    % A route with no stroke of its own never resolves a line width; one that
    % does must have resolved its own command's key, not a neighbour's.
    \ifx\KKSfitlwseen\relax\global\KKSfitusedlw=\z@\fi
    \setlength{\KKSfitdeclw}{\csname KKSfit@lw@#1\endcsname}%
    \ifdim\KKSfitusedlw>\dimexpr\KKSfitdeclw+\KKSfitslack\relax
      \KKSfitfail{#1/#3}{sized against a stroke of \the\KKSfitusedlw, but its
        outline is drawn with \the\KKSfitdeclw}%
    \fi
    \ifdim\KKSfitdeclw>\dimexpr\KKSfitusedlw+\KKSfitslack\relax
      \KKSfitfail{#1/#3}{sized against a stroke of \the\KKSfitusedlw, but its
        outline is drawn with \the\KKSfitdeclw}%
    \fi
    % Routes without a visual guard hand the reader the fitted box itself.
    \ifdim\KKSfitfinw<\z@
      \global\KKSfitfinw=\KKSfitgotw \global\KKSfitfinh=\KKSfitgoth
    \fi
    % (1) One \scalebox, hence one factor for both directions: the box that
    % comes out must have the aspect ratio the material went in with.
    \KKSfitsetratio{\KKSfitgotaspect}{\KKSfitgotw}{\KKSfitgoth}%
    \KKSfitsetratio{\KKSfitnataspect}{\KKSfitnatw}{\KKSfitnath}%
    \setlength{\KKSfitaspect}{%
      \dimexpr\KKSfitgotaspect-\KKSfitnataspect\relax}%
    \KKSfitabs{\KKSfitaspect}%
    \ifdim\KKSfitaspect>0.005\KKSfitnataspect
      \KKSfitfail{#1/#3}{aspect ratio changed from \strip@pt\KKSfitnataspect
        to \strip@pt\KKSfitgotaspect}%
    \fi
    % (2) The route may shrink, never enlarge past its own base size.
    \KKSfitsetratio{\KKSfitratio}{\KKSfitgotw}{\KKSfitnatw}%
    \setlength{\KKSfitshrunk}{\dimexpr#4pt-\KKSfitratio\relax}%
    \ifdim\KKSfitshrunk<-0.001pt
      \KKSfitfail{#1/#3}{scaled to \strip@pt\KKSfitratio, above the base #4}%
    \fi
    \ifdim\KKSfitratio<0.02pt
      \KKSfitfail{#1/#3}{scaled away to \strip@pt\KKSfitratio}%
    \fi
    % (3) The limit the command aimed at may never claim more room than the
    % outline that is actually drawn leaves free.
    \ifdim\dimexpr\KKSfitlima-\csname KKSfit@usable@#1\endcsname\relax
        >\KKSfitslack
      \KKSfitfail{#1/#3}{aimed at \the\KKSfitlima, but the drawn outline
        leaves only \csname KKSfit@usable@#1\endcsname}%
    \fi
    % (3b) Nor may it claim less.  Every setup is `the smaller of the
    % historical limit and what the outline leaves', and both of those are
    % measured here rather than taken from the package, so this is a two-sided
    % identity: a wrapper that resolves a neighbour's line width, or that
    % describes its corners by the wrong model, lands on the wrong side of it
    % whichever of the two widths is the larger.
    \pgfmathsetlength{\KKSfittmp}{%
      min(\csname KKSfit@cap@#1\endcsname, \csname KKSfit@usable@#1\endcsname)}%
    \ifdim\dimexpr\KKSfitlima-\KKSfittmp\relax>\KKSfitexactslack
      \KKSfitfail{#1/#3}{aimed at \the\KKSfitlima, but its own outline and
        inner box allow exactly \the\KKSfittmp}%
    \fi
    \ifdim\dimexpr\KKSfittmp-\KKSfitlima\relax>\KKSfitexactslack
      \KKSfitfail{#1/#3}{aimed at \the\KKSfitlima, but its own outline and
        inner box allow exactly \the\KKSfittmp}%
    \fi
    % (3c) The band is a shape constant, so it is checked the same way: the
    % circle allots the side of the square inscribed in its usable radius, and
    % every other shape allots nothing.  A setup that forgets to set the band,
    % or one that forgets to clear a neighbour's, fails here rather than
    % silently resizing that neighbour's material.
    \edef\KKSfitthiskind{\csname KKSfit@kind@#1\endcsname}%
    \def\KKSfitcirclekind{circle}%
    \ifx\KKSfitthiskind\KKSfitcirclekind
      \pgfmathsetlength{\KKSfittmp}{1.4142136*\KKSfitlima}%
      \ifdim\dimexpr\KKSfitlimband-\KKSfittmp\relax>\KKSfitexactslack
        \KKSfitfail{#1/#3}{banded at \the\KKSfitlimband, but the square
          inscribed in its usable circle has side \the\KKSfittmp}%
      \fi
      \ifdim\dimexpr\KKSfittmp-\KKSfitlimband\relax>\KKSfitexactslack
        \KKSfitfail{#1/#3}{banded at \the\KKSfitlimband, but the square
          inscribed in its usable circle has side \the\KKSfittmp}%
      \fi
      \ifdim\dimexpr\KKSfitlimwide-\KKSfitlimband\relax>\KKSfitexactslack
        \KKSfitfail{#1/#3}{takes \the\KKSfitlimwide\space as its wide width,
          but its band is \the\KKSfitlimband}%
      \fi
      \ifdim\dimexpr\KKSfitlimband-\KKSfitlimwide\relax>\KKSfitexactslack
        \KKSfitfail{#1/#3}{takes \the\KKSfitlimwide\space as its wide width,
          but its band is \the\KKSfitlimband}%
      \fi
    \else
      \ifdim\KKSfitlimband>\KKSfitslack
        \KKSfitfail{#1/#3}{banded at \the\KKSfitlimband, but this shape
          sets no band}%
      \fi
      \ifdim\KKSfitlimwide>\KKSfitslack
        \KKSfitfail{#1/#3}{takes \the\KKSfitlimwide\space as its wide width,
          but this shape has no wide width}%
      \fi
    \fi
    % (4) No overflow: the scaled box, centred on the shape, stays inside the
    % measured interior.
    \KKSfitjudgefit
    \ifnum\KKSfittaken=1
      \KKSfitsetreachaxis{\csname KKSfit@usable@#1\endcsname}%
    \else
      \KKSfitsetreach{\csname KKSfit@kind@#1\endcsname}%
        {\csname KKSfit@usable@#1\endcsname}{\csname KKSfit@cc@#1\endcsname}%
        {\csname KKSfit@rc@#1\endcsname}%
    \fi
    \ifdim\KKSfitreach>1.002pt
      \KKSfitfail{#1/#3}{reaches \strip@pt\KKSfitreach\space of the drawn interior
        (aimed at \the\KKSfitlima, the outline leaves
        \csname KKSfit@usable@#1\endcsname, and the box is
        \the\KKSfitgotw\space by \the\KKSfitgoth)}%
    \fi
    \edef\KKSfitouter{\strip@pt\KKSfitreach}
    % (4b) The same for the box the reader is actually handed.  On the routes
    % that carry a visual guard that box is not the one the fit produced: the
    % guard may take it in again after the fact.  It may only take it in --
    % never enlarge it, and never in one direction alone -- and the result
    % must still be inside the same measured interior, so that the check above
    % keeps its meaning for the finished symbol and not merely for the fit.
    \ifdim\dimexpr\KKSfitfinw-\KKSfitgotw\relax>\KKSfitslack
      \KKSfitfail{#1/#3}{the guard widened the fitted box from
        \the\KKSfitgotw\space to \the\KKSfitfinw}%
    \fi
    \ifdim\dimexpr\KKSfitfinh-\KKSfitgoth\relax>\KKSfitslack
      \KKSfitfail{#1/#3}{the guard heightened the fitted box from
        \the\KKSfitgoth\space to \the\KKSfitfinh}%
    \fi
    \KKSfitsetratio{\KKSfitfinaspect}{\KKSfitfinw}{\KKSfitfinh}%
    \setlength{\KKSfitaspect}{%
      \dimexpr\KKSfitfinaspect-\KKSfitgotaspect\relax}%
    \KKSfitabs{\KKSfitaspect}%
    \ifdim\KKSfitaspect>0.005\KKSfitgotaspect
      \KKSfitfail{#1/#3}{the guard changed the aspect ratio from
        \strip@pt\KKSfitgotaspect\space to \strip@pt\KKSfitfinaspect}%
    \fi
    \KKSfitjudgefinal
    \ifnum\KKSfittaken=1
      \KKSfitsetreachaxis{\csname KKSfit@usable@#1\endcsname}%
    \else
      \KKSfitsetreach{\csname KKSfit@kind@#1\endcsname}%
        {\csname KKSfit@usable@#1\endcsname}{\csname KKSfit@cc@#1\endcsname}%
        {\csname KKSfit@rc@#1\endcsname}%
    \fi
    \ifdim\KKSfitreach>1.002pt
      \KKSfitfail{#1/#3}{the published box reaches \strip@pt\KKSfitreach\space
        of the drawn interior}%
    \fi
    \edef\KKSfitpublic{\strip@pt\KKSfitreach}
    % (5) Nothing is thrown away either: material that had to be shrunk must
    % end up against the limit it was shrunk to, not somewhere below it.  This
    % judges the fit's own box: a guard that shrinks further does so on
    % purpose, and (4b) already holds it to shrinking only.
    \KKSfitjudgefit
    \ifnum\KKSfittaken=1
      \KKSfitsetreachaxis{\KKSfitlima}%
    \else
      \KKSfitsetreach{\csname KKSfit@kind@#1\endcsname}{\KKSfitlima}%
        {\KKSfitlimcc}{\KKSfitlimrc}%
    \fi
    % Material held by the band stops short of the region on purpose, and
    % then it is the band it has to fill instead.
    \setlength{\KKSfitbandreach}{\z@}%
    \ifdim\KKSfitlimband>\z@
      \KKSfitsetratio{\KKSfitbandreach}{\KKSfitgoth}{\KKSfitlimband}%
      \ifdim\KKSfitbandreach>1.002pt
        \KKSfitfail{#1/#3}{stands \strip@pt\KKSfitbandreach\space of the
          band \the\KKSfitlimband\space its shape allots}%
      \fi
    \fi
    % Material wider than its shape's wide width settles for the wide base;
    % that is the size it aimed at, so it is not holding back either.
    \KKSfitwideheldfalse
    \ifdim\KKSfitlimwide>\z@
      \ifdim\kks@fit@widebase\KKSfitnatw>\KKSfitlimwide\relax
        \ifdim\dimexpr\KKSfitratio-\kks@fit@widebase pt\relax<\KKSfitslack
          \ifdim\dimexpr\kks@fit@widebase pt-\KKSfitratio\relax<\KKSfitslack
            \KKSfitwideheldtrue
          \fi
        \fi
      \fi
    \fi
    \ifdim\KKSfitshrunk>0.001pt
      \ifdim\KKSfitreach<0.97pt
        \ifdim\KKSfitbandreach<0.97pt
          \ifKKSfitwideheld\else
            \KKSfitfail{#1/#3}{shrunk to \strip@pt\KKSfitratio, yet fills only
              \strip@pt\KKSfitreach\space of the room it aimed at and
              \strip@pt\KKSfitbandreach\space of its band}%
          \fi
        \fi
      \fi
    \fi
    \typeout{KKS-FIT:#1/#3 scale=\strip@pt\KKSfitratio\space
      interior=\KKSfitouter\space public=\KKSfitpublic\space
      limit=\strip@pt\KKSfitreach\space
      band=\strip@pt\KKSfitbandreach}%
  \endgroup}

% ---------------------------------------------------------------------------
% The enclosures, with the geometry their drawing nodes state.
% ---------------------------------------------------------------------------
% The line width the options really resolve to, so that a fixture which sets
% `linewidth' or a per-command key measures the outline it asked for.
\newcommand{\KKSfitlw}[2]{%
  \kks@linewidth{\csname kksymbols@#1\endcsname}{\kksymbols@linewidth}{#2}}
\newcommand{\KKSfitshapes}{%
  \KKSfitshape{maru}{circle}{shape=circle,%
      line width=\KKSfitlw{marulinewidth}{.04\zw},minimum size=.9\zw,%
      inner sep=.02\zw}{0.87\zw/sqrt(2)}{1}{0pt}%
    {\KKSfitlw{marulinewidth}{.04\zw}}%
  \KKSfitshape{kuromaru}{circle}{shape=circle,%
      line width=\KKSfitlw{kuromarulinewidth}{.04\zw},minimum size=.9\zw,%
      inner sep=.02\zw}{0.87\zw/sqrt(2)}{1}{0pt}%
    {\KKSfitlw{kuromarulinewidth}{.04\zw}}%
  \KKSfitshape{nmaru}{circle}{shape=circle,%
      line width=\KKSfitlw{nmarulinewidth}{.03\zw},minimum size=.88\zw,%
      double,double distance=.03\zw,inner sep=.06\zw,scale=.9}%
    {0.82\zw/sqrt(2)}{0.9}{0pt}{\KKSfitlw{nmarulinewidth}{.03\zw}}%
  \KKSfitshape{seihou}{rect}{line width=\KKSfitlw{seihoulinewidth}{.04\zw},%
      minimum size=.85\zw,inner sep=.1\zw}{0.68\zw}{1}{0pt}%
    {\KKSfitlw{seihoulinewidth}{.04\zw}}%
  \KKSfitshape{kuroseihou}{rect}{%
      line width=\KKSfitlw{kuroseihoulinewidth}{.04\zw},minimum size=.85\zw,%
      inner sep=.1\zw}{0.68\zw}{1}{0pt}%
    {\KKSfitlw{kuroseihoulinewidth}{.04\zw}}%
  \KKSfitshape{seimaru}{rect}{line width=\KKSfitlw{seimarulinewidth}{.04\zw},%
      minimum size=.85\zw,inner sep=.1\zw,rounded corners=.13\zw}%
    {0.68\zw}{1}{.13\zw}{\KKSfitlw{seimarulinewidth}{.04\zw}}%
  \KKSfitshape{kuroseimaru}{rect}{%
      line width=\KKSfitlw{kuroseimarulinewidth}{.04\zw},minimum size=.85\zw,%
      inner sep=.1\zw,rounded corners=.13\zw}{0.68\zw}{1}{.13\zw}%
    {\KKSfitlw{kuroseimarulinewidth}{.04\zw}}%
  \KKSfitshape{hishi}{diamond}{shape=diamond,%
      line width=\KKSfitlw{hishilinewidth}{.04\zw},minimum size=.5\zw,%
      inner sep=.03\zw}{0.85\zw/2}{1}{0pt}%
    {\KKSfitlw{hishilinewidth}{.04\zw}}%
  \KKSfitshape{kurohishi}{diamond}{shape=diamond,%
      line width=\KKSfitlw{kurohishilinewidth}{.04\zw},minimum size=.5\zw,%
      inner sep=.03\zw}{0.85\zw/2}{1}{0pt}%
    {\KKSfitlw{kurohishilinewidth}{.04\zw}}%
  \KKSfitshape{maruhishi}{diamond}{shape=diamond,%
      line width=\KKSfitlw{maruhishilinewidth}{.04\zw},minimum size=.5\zw,%
      inner sep=.03\zw}{0.85\zw/2}{1}{.06\zw}%
    {\KKSfitlw{maruhishilinewidth}{.04\zw}}%
  \KKSfitshape{kuromaruhishi}{diamond}{shape=diamond,%
      line width=\KKSfitlw{kuromaruhishilinewidth}{.04\zw},%
      minimum size=.5\zw,inner sep=.03\zw}{0.85\zw/2}{1}{.06\zw}%
    {\KKSfitlw{kuromaruhishilinewidth}{.04\zw}}%
  % The brackets and \ichimoji draw no frame, so their opening is exactly the
  % box the enclosure reserves: a stroke-free rectangle of that size.
  \KKSfitshape{kakko}{rect}{line width=0pt,minimum size=0pt,inner sep=0pt}%
    {0.8\zw}{1}{0pt}{0pt}%
  \KKSfitshape{ichimoji}{rect}{line width=0pt,minimum size=0pt,inner sep=0pt}%
    {\zw}{1}{0pt}{0pt}%
}

% Two enclosures really are drawn with different strokes.  A fixture that sets
% out to tell one command's line-width key from another's is worthless if the
% widths it chose happen to land on the same outline, so it says so here and
% the suite refuses to pretend otherwise.
\newcommand{\KKSfitdistinct}[2]{%
  \begingroup
    \pgfmathsetlength{\KKSfittmp}{%
      abs(\csname KKSfit@stroke@#1\endcsname
          - \csname KKSfit@stroke@#2\endcsname)}%
    \ifdim\KKSfittmp<0.15pt
      \KKSfitfail{#1/#2}{are drawn with strokes only \the\KKSfittmp\space
        apart, which is too little to tell their keys apart}%
    \fi
  \endgroup}
% A corner arc of radius #2 would really be visible on enclosure #1.  A tip
% rounded on an arc the pen is wider than is a tip the pen has already
% squared off: the two corner models then genuinely coincide, and a fixture
% that meant to tell them apart has to be drawn with a thinner stroke.  The
% .414 is the share of the arc a diamond's tip loses along its diagonal.
\newcommand{\KKSfitcorneractive}[2]{%
  \begingroup
    \pgfmathsetlength{\KKSfittmp}{%
      0.4142136*max(0, (#2) - .5*\csname KKSfit@stroke@#1\endcsname)}%
    \ifdim\KKSfittmp<0.05pt
      \KKSfitfail{#1}{a corner radius of #2 would move its interior by only
        \the\KKSfittmp\space against the stroke it is drawn with, which is
        too little to tell a rounded tip from a sharp one}%
    \fi
  \endgroup}

% ---------------------------------------------------------------------------
% The suite.  \KKSfitrow runs one piece of material through every star route
% whose star selects the safe-sizing path; \jegg is deliberately absent
% because its star selects a gray background, not a size.
% ---------------------------------------------------------------------------
\newcommand{\KKSfitrow}[1]{%
  \KKSfitcheck{maru}{\maru*}{#1}{.9}%
  \KKSfitcheck{kuromaru}{\kuromaru*}{#1}{.9}%
  \KKSfitcheck{nmaru}{\nmaru*}{#1}{.9}%
  \KKSfitcheck{seihou}{\seihou*}{#1}{.9}%
  \KKSfitcheck{kuroseihou}{\kuroseihou*}{#1}{.9}%
  \KKSfitcheck{seimaru}{\seimaru*}{#1}{.9}%
  \KKSfitcheck{kuroseimaru}{\kuroseimaru*}{#1}{.9}%
  \KKSfitcheck{hishi}{\hishi*}{#1}{.9}%
  \KKSfitcheck{kurohishi}{\kurohishi*}{#1}{.9}%
  \KKSfitcheck{maruhishi}{\maruhishi*}{#1}{.9}%
  \KKSfitcheck{kuromaruhishi}{\kuromaruhishi*}{#1}{.9}%
  \KKSfitcheck{kakko}{\kakko*}{#1}{.9}%
  \KKSfitcheck{kakko}{\sumikakko*}{#1}{.9}%
  \KKSfitcheck{kakko}{\kakukakko*}{#1}{.9}%
  \KKSfitcheck{kakko}{\kikakko*}{#1}{.9}%
  \KKSfitcheck{kakko}{\ykakko*}{#1}{.9}%
  \KKSfitcheck{kakko}{\nykakko*}{#1}{.9}%
  \KKSfitcheck{kakko}{\namikakko*}{#1}{.9}%
  \KKSfitcheck{kakko}{\kagikakko*}{#1}{.9}%
  \KKSfitcheck{kakko}{\nkagikakko*}{#1}{.9}%
  \KKSfitcheck{ichimoji}{\ichimoji*}{#1}{1}%
}
\newcommand{\KKSfitsuite}{%
  \begingroup
    \KKSfitshapes
    % The reported case, then a wider string, a tall one, one that is both,
    % a wide single letter, a narrow one, and a full-width character.
    \KKSfitrow{abc}%
    \KKSfitrow{mmmm}%
    \KKSfitrow{Qjgy}%
    \KKSfitrow{Wgabcdef}%
    \KKSfitrow{m}%
    \KKSfitrow{i}%
    \KKSfitrow{漢}%
    % The three letters called out as critical, each on its own and then in
    % combinations that put their different metrics against one another: `c'
    % is short and sits wholly on the baseline, `m' is the widest of the three
    % and equally short, and `j' is the only one with a descender and with ink
    % above the x-height, so a route that measured height without depth would
    % size `j' as if it were `c'.  The height recorded throughout is the box's
    % height plus its depth, which is what makes these rows bite.
    \KKSfitrow{c}%
    \KKSfitrow{j}%
    \KKSfitrow{cmj}%
    \KKSfitrow{mjm}%
    \KKSfitrow{cjc}%
    \KKSfitrow{jjj}%
    \typeout{KKS-STAR-FIT:OK}%
  \endgroup}

% With a stroke heavy enough to swallow an interior there is nothing left to
% fit into, so the routes must fall back to their plain base size instead of
% dividing by a non-positive limit.  Only the properties that still make sense
% are asserted: the scale is the base one, and it is still uniform.
\newcommand{\KKSfitbasecheck}[4]{%
  \begingroup
    \global\let\KKSfitseen\relax
    \setbox\z@=\hbox{#2{#3}}%
    \ifx\KKSfitseen\relax
      \KKSfitfail{#1/#3}{the star route never reached the uniform fit}%
    \fi
    \KKSfitsetratio{\KKSfitgotaspect}{\KKSfitgotw}{\KKSfitgoth}%
    \KKSfitsetratio{\KKSfitnataspect}{\KKSfitnatw}{\KKSfitnath}%
    \setlength{\KKSfitaspect}{%
      \dimexpr\KKSfitgotaspect-\KKSfitnataspect\relax}%
    \KKSfitabs{\KKSfitaspect}%
    \ifdim\KKSfitaspect>0.005\KKSfitnataspect
      \KKSfitfail{#1/#3}{aspect ratio changed from \strip@pt\KKSfitnataspect
        to \strip@pt\KKSfitgotaspect}%
    \fi
    \KKSfitsetratio{\KKSfitratio}{\KKSfitgotw}{\KKSfitnatw}%
    \setlength{\KKSfitshrunk}{\dimexpr\KKSfitratio-#4pt\relax}%
    \KKSfitabs{\KKSfitshrunk}%
    \ifdim\KKSfitshrunk>0.001pt
      \KKSfitfail{#1/#3}{fell back to \strip@pt\KKSfitratio, not to the base
        size #4}%
    \fi
  \endgroup}
\newcommand{\KKSfitbaserow}[1]{%
  \KKSfitbasecheck{maru}{\maru*}{#1}{.9}%
  \KKSfitbasecheck{kuromaru}{\kuromaru*}{#1}{.9}%
  \KKSfitbasecheck{nmaru}{\nmaru*}{#1}{.9}%
  \KKSfitbasecheck{seihou}{\seihou*}{#1}{.9}%
  \KKSfitbasecheck{kuroseihou}{\kuroseihou*}{#1}{.9}%
  \KKSfitbasecheck{seimaru}{\seimaru*}{#1}{.9}%
  \KKSfitbasecheck{kuroseimaru}{\kuroseimaru*}{#1}{.9}%
  \KKSfitbasecheck{hishi}{\hishi*}{#1}{.9}%
  \KKSfitbasecheck{kurohishi}{\kurohishi*}{#1}{.9}%
  \KKSfitbasecheck{maruhishi}{\maruhishi*}{#1}{.9}%
  \KKSfitbasecheck{kuromaruhishi}{\kuromaruhishi*}{#1}{.9}%
}
\newcommand{\KKSfitbasesuite}{%
  \begingroup
    \KKSfitbaserow{abc}%
    \KKSfitbaserow{c}%
    \KKSfitbaserow{m}%
    \KKSfitbaserow{j}%
    \KKSfitbaserow{cmj}%
    \KKSfitbaserow{漢}%
    \typeout{KKS-STAR-FIT-GUARD:OK}%
  \endgroup}

% Visual counterpart: one column per enclosure, in the order of the bug
% report, so that a rendered page can be compared with it directly.
\newcommand{\KKSfitshowrow}[1]{%
  \fbox{\maru*{#1}}\fbox{\kuromaru*{#1}}\fbox{\nmaru*{#1}}\fbox{\jegg*{#1}}%
  \fbox{\seihou*{#1}}\fbox{\kuroseihou*{#1}}%
  \fbox{\seimaru*{#1}}\fbox{\kuroseimaru*{#1}}%
  \fbox{\hishi*{#1}}\fbox{\kurohishi*{#1}}%
  \fbox{\maruhishi*{#1}}\fbox{\kuromaruhishi*{#1}}%
  \fbox{\kakko*{#1}}\fbox{\ichimoji*{#1}}}
\newcommand{\KKSfitshowcolumn}[1]{%
  \fbox{\maru*{#1}}\par\fbox{\kuromaru*{#1}}\par\fbox{\nmaru*{#1}}\par
  \fbox{\jegg*{#1}}\par\fbox{\seihou*{#1}}\par\fbox{\kuroseihou*{#1}}\par
  \fbox{\seimaru*{#1}}\par\fbox{\kuroseimaru*{#1}}\par
  \fbox{\hishi*{#1}}\par\fbox{\kurohishi*{#1}}\par
  \fbox{\maruhishi*{#1}}\par\fbox{\kuromaruhishi*{#1}}\par}
\makeatother
